\part{III--IV klasė}

\section{Natūralieji skaičiai iki milijono}\label{sec:naturalieji-skaiciai-iki-milijono}

\subsection{Apibrėžtys}

\begin{apibreztis}
Natūralieji skaičiai yra skaičiai, kuriais skaičiuojame daiktus arba jų kiekį:
\[
1,2,3,4,\ldots
\]
III--IV klasėje mokomasi tvirtai dirbti su skaičiais iki \(10\,000\), o vėliau pažinti ir skaičius iki \(1\,000\,000\).
\end{apibreztis}

\begin{apibreztis}
Skaitmens vietinė vertė rodo, kiek skaitmuo reiškia pagal savo vietą skaičiuje. Skaičiuje \(482\,615\) skaitmuo \(4\) reiškia \(400\,000\), skaitmuo \(8\) reiškia \(80\,000\), skaitmuo \(2\) reiškia \(2\,000\), skaitmuo \(6\) reiškia \(600\), skaitmuo \(1\) reiškia \(10\), o skaitmuo \(5\) reiškia \(5\).
\end{apibreztis}

\begin{apibreztis}
Skaičiaus skyrių suma yra skaičiaus užrašymas sudedant jo vietines vertes, pavyzdžiui
\[
306\,048=300\,000+6\,000+40+8.
\]
\end{apibreztis}

\begin{pastaba}
Didelius skaičius patogu skaidyti grupėmis po tris skaitmenis nuo dešinės: \(724381\) rašome \(724\,381\). Tarpas čia nėra veiksmo ženklas, jis tik padeda skaityti skaičių.
\end{pastaba}

\subsection{Teoremos ir savybės}

\begin{savybe}
Jei vienas natūralusis skaičius turi daugiau skaitmenų negu kitas, jis yra didesnis. Pavyzdžiui, \(99\,999<100\,000\).
\end{savybe}

\begin{savybe}
Jei du skaičiai turi tiek pat skaitmenų, juos lyginame iš kairės į dešinę. Pirmas nevienodas skaitmuo parodo, kuris skaičius didesnis.
\end{savybe}

\begin{savybe}
Skaičių galima keisti iš skaitmenų užrašo į skyrių sumą ir atgal. Tai padeda suprasti, kodėl veikia sudėties, atimties, apvalinimo ir stulpelinio skaičiavimo būdai.
\end{savybe}

\subsection{Taisyklės ir algoritmai}

\textbf{Skaičiaus skaitymas.}
\begin{enumerate}[label=\arabic*.]
\item Suskirstyk skaičių į grupes po tris skaitmenis nuo dešinės.
\item Kairiąją grupę skaityk kaip tūkstančių arba milijonų grupę.
\item Dešiniąją grupę skaityk kaip vienetų grupę.
\item Jei grupėje yra nulis, jo atskirai neskaityk, bet jo vietinė vertė vis tiek svarbi.
\end{enumerate}

\textbf{Skaičių palyginimas.}
\begin{enumerate}[label=\arabic*.]
\item Palygink skaitmenų skaičių.
\item Jei skaitmenų tiek pat, lygink iš kairės: šimtus tūkstančių, dešimtis tūkstančių, tūkstančius, šimtus, dešimtis, vienetus.
\item Kai randi pirmą skirtingą skaitmenį, nuspręsk, kuris skaičius didesnis.
\end{enumerate}

\textbf{Apvalinimas.}
Norėdamas apvalinti iki dešimčių, šimtų, tūkstančių ar didesnio skyriaus, žiūrėk į skaitmenį, esantį iš karto dešinėje nuo apvalinamos vietos. Jei jis yra \(0,1,2,3,4\), apvalinamas skaitmuo nesikeičia. Jei jis yra \(5,6,7,8,9\), apvalinamas skaitmuo padidėja vienetu. Visi dešinėje esantys skaitmenys tampa nuliais.

\subsection{Iliustruojantys pavyzdžiai}

\textbf{1 pavyzdys.} Išskaidykime skaičių \(438\,205\) skyrių suma.

Skaitmuo \(4\) yra šimtų tūkstančių vietoje, todėl reiškia \(400\,000\). Skaitmuo \(3\) reiškia \(30\,000\), skaitmuo \(8\) reiškia \(8\,000\), skaitmuo \(2\) reiškia \(200\), skaitmuo \(0\) dešimčių neprideda, o skaitmuo \(5\) reiškia \(5\). Todėl
\[
438\,205=400\,000+30\,000+8\,000+200+5.
\]
\textbf{Atsakymas:} \(400\,000+30\,000+8\,000+200+5\).

\textbf{2 pavyzdys.} Palyginkime \(572\,419\) ir \(572\,491\).

Abu skaičiai turi po šešis skaitmenis. Lyginame iš kairės:
\[
5=5,\quad 7=7,\quad 2=2,\quad 4=4.
\]
Toliau lyginame dešimčių skaitmenis: pirmame skaičiuje yra \(1\), antrame \(9\). Kadangi \(1<9\), tai
\[
572\,419<572\,491.
\]
\textbf{Atsakymas:} \(572\,419<572\,491\).

\textbf{3 pavyzdys.} Suapvalinkime \(286\,735\) iki tūkstančių.

Tūkstančių skaitmuo yra \(6\). Dešinėje nuo jo yra šimtų skaitmuo \(7\). Kadangi \(7\geq 5\), tūkstančių skaitmenį padidiname vienetu: \(6\) tampa \(7\). Visi dešinieji skaitmenys tampa nuliais:
\[
286\,735\approx 287\,000.
\]
\textbf{Atsakymas:} \(287\,000\).

\subsection{Dažnos klaidos ir savikontrolė}

\textbf{Dažnos klaidos.}
\begin{enumerate}[label=\arabic*.]
\item Praleidžiamas nulis skaičiaus viduryje: \(405\,018\) nėra tas pats kaip \(45\,018\).
\item Lyginama tik paskutiniai skaitmenys, nors pirmiausia reikia lyginti iš kairės.
\item Apvalinant pamirštama dešinėje esančius skaitmenis pakeisti nuliais.
\end{enumerate}

\textbf{Pasitikrinkite.}
\begin{enumerate}[label=\arabic*.]
\item Užrašykite skyrių suma: \(709\,034\).
\item Įrašykite ženklą \(<\), \(>\) arba \(=\): \(84\,902\ \square\ 84\,920\).
\item Suapvalinkite \(53\,681\) iki šimtų.
\end{enumerate}

\subsection{Ryšys su kitomis temomis}

Ši tema remiasi skaičiais iki \(1000\) ir vietine verte iš skyriaus \ref{sec:skaiciai-iki-1000-ir-vietine-verte}. Ji reikalinga veiksmams skyriuje \ref{sec:aritmetiniai-veiksmai-ir-savybes}, stulpelinei daugybai ir dalybai skyriuje \ref{sec:daugyba-ir-dalyba-stulpeliu}, finansiniams skaičiavimams skyriuje \ref{sec:finansiniai-skaiciavimai-kainos-pokyciai-sprendimai} ir duomenų skaitymui skyriuje \ref{sec:duomenu-lenteles-ir-diagramos}.

\section{Aritmetiniai veiksmai ir savybės}\label{sec:aritmetiniai-veiksmai-ir-savybes}

\subsection{Apibrėžtys}

\begin{apibreztis}
Sudėtis yra veiksmas, kai kelis kiekius sujungiame į vieną. Sudėties rezultatas vadinamas suma, o sudedami skaičiai -- dėmenimis.
\end{apibreztis}

\begin{apibreztis}
Atimtis yra veiksmas, kai iš kiekio atimame dalį arba randame dviejų kiekių skirtumą. Atimties rezultatas vadinamas skirtumu.
\end{apibreztis}

\begin{apibreztis}
Daugyba yra vienodų dėmenų sudėtis. Dalyba yra veiksmas, kai kiekį dalijame į lygias dalis arba nustatome, kiek kartų vienas skaičius telpa kitame.
\end{apibreztis}

\begin{apibreztis}
Veiksmų tvarka nurodo, kuriuos veiksmus reiškinyje atliekame pirmiau.
\end{apibreztis}

\subsection{Teoremos ir savybės}

\begin{savybe}
Sudėties perstatomumo savybė:
\[
a+b=b+a.
\]
Pavyzdžiui, \(36+24=24+36\).
\end{savybe}

\begin{savybe}
Sudėties jungiamumo savybė:
\[
(a+b)+c=a+(b+c).
\]
Ji leidžia patogiai grupuoti dėmenis.
\end{savybe}

\begin{savybe}
Daugybos perstatomumo savybė:
\[
a\cdot b=b\cdot a.
\]
Pavyzdžiui, \(7\cdot 8=8\cdot 7\).
\end{savybe}

\begin{savybe}
Daugybos skirstomumo sudėties atžvilgiu savybė:
\[
a\cdot(b+c)=a\cdot b+a\cdot c.
\]
Ji padeda skaičiuoti mintinai ir pasiruošti daugybai stulpeliu.
\end{savybe}

\subsection{Taisyklės ir algoritmai}

\textbf{Veiksmų tvarka.}
\begin{enumerate}[label=\arabic*.]
\item Pirmiausia atliekami veiksmai skliaustuose.
\item Tada atliekama daugyba ir dalyba iš kairės į dešinę.
\item Galiausiai atliekama sudėtis ir atimtis iš kairės į dešinę.
\end{enumerate}

\textbf{Tikrinimas atvirkštiniu veiksmu.}
Sudėtį tikriname atimtimi, atimtį -- sudėtimi, daugybą -- dalyba, dalybą -- daugyba. Jei \(48+35=83\), tikriname \(83-35=48\). Jei \(72:8=9\), tikriname \(9\cdot8=72\).

\subsection{Iliustruojantys pavyzdžiai}

\textbf{1 pavyzdys.} Apskaičiuokime \(48+27+52\) patogiai.

Pastebime, kad \(48+52=100\). Todėl dėmenis galime perstatyti ir sugrupuoti:
\[
48+27+52=(48+52)+27=100+27=127.
\]
\textbf{Atsakymas:} \(127\).

\textbf{2 pavyzdys.} Apskaičiuokime \(6\cdot 23\) naudodami skirstomumo savybę.

Skaičių \(23\) išskaidome į \(20+3\):
\[
6\cdot23=6\cdot(20+3)=6\cdot20+6\cdot3=120+18=138.
\]
\textbf{Atsakymas:} \(138\).

\textbf{3 pavyzdys.} Apskaičiuokime \(80-24:6\).

Pirmiausia atliekame dalybą:
\[
24:6=4.
\]
Tada atimame:
\[
80-4=76.
\]
\textbf{Atsakymas:} \(76\).

\textbf{4 pavyzdys.} Klasė pirko \(4\) sąsiuvinių rinkinius po \(8\) sąsiuvinius ir dar \(6\) sąsiuvinius. Kiek sąsiuvinių nupirkta?

Pirmiausia randame sąsiuvinius rinkiniuose:
\[
4\cdot8=32.
\]
Pridedame dar \(6\):
\[
32+6=38.
\]
\textbf{Atsakymas:} \(38\) sąsiuviniai.

\subsection{Dažnos klaidos ir savikontrolė}

\textbf{Dažnos klaidos.}
\begin{enumerate}[label=\arabic*.]
\item Reiškinyje \(80-24:6\) pirma atliekama atimtis, nors turi būti atliekama dalyba.
\item Sumaišomi atvirkštiniai veiksmai: dalyba tikrinama ne atimtimi, o daugyba.
\item Skirstomumo savybėje padauginama tik viena suma dalis: \(6\cdot(20+3)\neq 6\cdot20+3\).
\end{enumerate}

\textbf{Pasitikrinkite.}
\begin{enumerate}[label=\arabic*.]
\item Apskaičiuokite patogiai: \(25+39+75\).
\item Apskaičiuokite: \(9\cdot(30+4)\).
\item Apskaičiuokite: \(100-36:4\).
\end{enumerate}

\subsection{Ryšys su kitomis temomis}

Aritmetiniai veiksmai naudojami beveik visur: skaičių iki milijono temoje \ref{sec:naturalieji-skaiciai-iki-milijono}, daugyboje ir dalyboje stulpeliu \ref{sec:daugyba-ir-dalyba-stulpeliu}, reiškiniuose \ref{sec:reiskiniai-ir-raidiniai-zymejimai}, matavimuose \ref{sec:perimetras-plotas-turis-ir-masto-pradzia} ir duomenų uždaviniuose \ref{sec:duomenu-lenteles-ir-diagramos}.

\section{Daugyba ir dalyba stulpeliu}\label{sec:daugyba-ir-dalyba-stulpeliu}

\subsection{Apibrėžtys}

\begin{apibreztis}
Daugyba stulpeliu yra rašytinis būdas dauginti didesnius skaičius, kai skaičiuojama pagal vietines vertes: vienetus, dešimtis, šimtus ir taip toliau.
\end{apibreztis}

\begin{apibreztis}
Dalyba stulpeliu yra rašytinis būdas dalyti didesnius skaičius, kai dalijame nuo didžiausios vietinės vertės, dauginame, atimame ir nuleidžiame kitą skaitmenį.
\end{apibreztis}

\begin{apibreztis}
Liekana yra skaičius, kuris lieka padalijus, kai jo nebegalima padalyti į sveikas lygias dalis. Liekana visada mažesnė už daliklį.
\end{apibreztis}

\subsection{Teoremos ir savybės}

\begin{savybe}
Dauginant skaičių iš \(10\), \(100\), \(1000\), prie skaičiaus dešinės prirašome atitinkamai vieną, du arba tris nulius:
\[
37\cdot100=3\,700.
\]
\end{savybe}

\begin{savybe}
Dalybos su liekana tikrinimo taisyklė:
\[
\text{daliklis}\cdot\text{dalmuo}+\text{liekana}=\text{dalinys}.
\]
\end{savybe}

\begin{savybe}
Dauginant iš dviženklio skaičiaus, gaunami daliniai rezultatai pagal vienetus ir dešimtis. Juos būtina sudėti atsižvelgiant į vietinę vertę.
\end{savybe}

\subsection{Taisyklės ir algoritmai}

\textbf{Daugyba stulpeliu iš vienaženklio skaičiaus.}
\begin{enumerate}[label=\arabic*.]
\item Pradėk nuo vienetų.
\item Padaugink skaitmenį iš daugiklio.
\item Vienetus rašyk, dešimtis kelk į kitą skyrių.
\item Tęsk su dešimtimis, šimtais ir kitais skyriais.
\end{enumerate}

\textbf{Daugyba iš dviženklio skaičiaus.}
\begin{enumerate}[label=\arabic*.]
\item Padaugink iš vienetų skaitmens.
\item Padaugink iš dešimčių skaitmens, prisimindamas, kad tai dešimtys.
\item Sudėk dalinius rezultatus.
\end{enumerate}

\textbf{Dalyba stulpeliu.}
\begin{enumerate}[label=\arabic*.]
\item Imk tiek pirmųjų dalinio skaitmenų, kad būtų galima dalyti iš daliklio.
\item Parašyk dalmens skaitmenį.
\item Padaugink daliklį iš parašyto skaitmens ir atimk.
\item Nuleisk kitą dalinio skaitmenį.
\item Kartok, kol neliks skaitmenų.
\item Patikrink atsakymą.
\end{enumerate}

\subsection{Iliustruojantys pavyzdžiai}

\textbf{1 pavyzdys.} Sudauginkime \(236\cdot4\).

Dauginame nuo vienetų:
\[
6\cdot4=24,
\]
rašome \(4\), keliame \(2\). Toliau
\[
3\cdot4+2=12+2=14,
\]
rašome \(4\), keliame \(1\). Toliau
\[
2\cdot4+1=9.
\]
Taigi
\[
236\cdot4=944.
\]
\textbf{Atsakymas:} \(944\).

\textbf{2 pavyzdys.} Sudauginkime \(128\cdot23\).

Pirma dauginame iš \(3\):
\[
128\cdot3=384.
\]
Tada dauginame iš \(20\):
\[
128\cdot20=2\,560.
\]
Sudedame:
\[
384+2\,560=2\,944.
\]
\textbf{Atsakymas:} \(2\,944\).

\textbf{3 pavyzdys.} Padalykime \(845:6\) su liekana.

Imame \(8\): \(8:6=1\), lieka \(2\). Nuleidžiame \(4\), gauname \(24\). \(24:6=4\), lieka \(0\). Nuleidžiame \(5\). \(5:6=0\), todėl dalmens vienetų vietoje rašome \(0\), lieka \(5\). Gauname
\[
845:6=140\ \text{liek.}\ 5.
\]
Tikriname:
\[
6\cdot140+5=840+5=845.
\]
\textbf{Atsakymas:} \(140\), liekana \(5\).

\subsection{Dažnos klaidos ir savikontrolė}

\textbf{Dažnos klaidos.}
\begin{enumerate}[label=\arabic*.]
\item Dauginant iš dviženklio skaičiaus pamirštama, kad dauginimas iš dešimčių reiškia dauginimą iš \(10\) kartų didesnio skaičiaus.
\item Dalyboje neparašomas nulis dalmens viduryje arba gale.
\item Liekana paliekama didesnė arba lygi dalikliui, nors taip būti negali.
\end{enumerate}

\textbf{Pasitikrinkite.}
\begin{enumerate}[label=\arabic*.]
\item Apskaičiuokite: \(347\cdot6\).
\item Apskaičiuokite: \(214\cdot32\).
\item Padalykite su liekana: \(739:8\).
\end{enumerate}

\subsection{Ryšys su kitomis temomis}

Daugyba ir dalyba stulpeliu remiasi vietine verte iš skyriaus \ref{sec:naturalieji-skaiciai-iki-milijono} ir veiksmų savybėmis iš skyriaus \ref{sec:aritmetiniai-veiksmai-ir-savybes}. Šie įgūdžiai reikalingi trupmenoms \ref{sec:trupmenos-ir-desimtaines-trupmenos-pradzia}, matavimams \ref{sec:matavimo-vienetai-laikas-greitis-plotas-turis} ir vėlesniam daliklių bei kartotinių mokymuisi.

\section{Trupmenos ir dešimtainės trupmenos pradžia}\label{sec:trupmenos-ir-desimtaines-trupmenos-pradzia}

\subsection{Apibrėžtys}

\begin{apibreztis}
Trupmena rodo, į kiek lygių dalių padalytas visas daiktas ar kiekis ir kiek tokių dalių paimta. Trupmenoje \(\frac{3}{4}\) skaičius \(4\) yra vardiklis, o skaičius \(3\) yra skaitiklis.
\end{apibreztis}

\begin{apibreztis}
Vardiklis parodo, į kiek lygių dalių padalyta visuma. Skaitiklis parodo, kiek tokių dalių paimta.
\end{apibreztis}

\begin{apibreztis}
Dešimtainė trupmena yra trupmenos rašymas naudojant kablelį. Pavyzdžiui, \(0,5\) reiškia penkias dešimtąsias dalis, tai yra \(\frac{5}{10}\).
\end{apibreztis}

\begin{apibreztis}
Mišrusis skaičius turi sveikąją dalį ir trupmeninę dalį, pavyzdžiui \(2\frac{1}{3}\).
\end{apibreztis}

\subsection{Teoremos ir savybės}

\begin{savybe}
Jei dviejų trupmenų vardikliai vienodi, didesnė ta trupmena, kurios skaitiklis didesnis:
\[
\frac{3}{8}>\frac{2}{8}.
\]
\end{savybe}

\begin{savybe}
Viena antroji yra tiek pat, kiek penkios dešimtosios:
\[
\frac{1}{2}=\frac{5}{10}=0,5.
\]
\end{savybe}

\begin{savybe}
Trupmenos su vardikliu \(10\) tiesiogiai užrašomos dešimtainėmis trupmenomis:
\[
\frac{1}{10}=0,1,\quad \frac{7}{10}=0,7,\quad 3\frac{4}{10}=3,4.
\]
\end{savybe}

\subsection{Taisyklės ir algoritmai}

\textbf{Visumos dalies radimas.}
\begin{enumerate}[label=\arabic*.]
\item Visumą padalyk iš vardiklio.
\item Gautą vienos dalies dydį padaugink iš skaitiklio.
\end{enumerate}

\textbf{Trupmenų su vienodais vardikliais sudėtis.}
Sudedame skaitiklius, o vardiklį paliekame tą patį:
\[
\frac{a}{n}+\frac{b}{n}=\frac{a+b}{n}.
\]

\textbf{Dešimtųjų skaitymas.}
Skaičius \(4,6\) reiškia \(4\) sveikuosius ir \(6\) dešimtąsias dalis. Jį galima užrašyti
\[
4,6=4\frac{6}{10}.
\]

\subsection{Iliustruojantys pavyzdžiai}

\textbf{1 pavyzdys.} Raskime \(\frac{3}{5}\) skaičiaus \(20\).

Pirmiausia randame vieną penktadalį:
\[
20:5=4.
\]
Tada paimame tris tokias dalis:
\[
4\cdot3=12.
\]
\textbf{Atsakymas:} \(12\).

\textbf{2 pavyzdys.} Sudėkime \(\frac{2}{7}+\frac{3}{7}\).

Vardikliai vienodi, todėl sudedame tik skaitiklius:
\[
\frac{2}{7}+\frac{3}{7}=\frac{2+3}{7}=\frac{5}{7}.
\]
\textbf{Atsakymas:} \(\frac{5}{7}\).

\textbf{3 pavyzdys.} Užrašykime \(4\frac{6}{10}\) dešimtaine trupmena.

Sveikoji dalis yra \(4\). Trupmeninė dalis \(\frac{6}{10}\) reiškia \(0,6\). Todėl
\[
4\frac{6}{10}=4,6.
\]
\textbf{Atsakymas:} \(4,6\).

\textbf{4 pavyzdys.} Kuri dalis didesnė: \(\frac{5}{10}\) ar \(\frac{8}{10}\)?

Vardikliai vienodi, todėl lyginame skaitiklius. Kadangi \(5<8\), tai
\[
\frac{5}{10}<\frac{8}{10}.
\]
\textbf{Atsakymas:} \(\frac{8}{10}\) yra didesnė.

\subsection{Dažnos klaidos ir savikontrolė}

\textbf{Dažnos klaidos.}
\begin{enumerate}[label=\arabic*.]
\item Vardiklis suprantamas kaip paimtų dalių skaičius, nors jis rodo visų lygių dalių skaičių.
\item Sudedant vienodus vardiklius klaidingai sudedami ir vardikliai: \(\frac{2}{7}+\frac{3}{7}\neq\frac{5}{14}\).
\item Dešimtainėje trupmenoje kablelis painiojamas su tūkstančių skirtuku.
\end{enumerate}

\textbf{Pasitikrinkite.}
\begin{enumerate}[label=\arabic*.]
\item Raskite \(\frac{2}{3}\) skaičiaus \(18\).
\item Apskaičiuokite: \(\frac{4}{9}+\frac{2}{9}\).
\item Užrašykite dešimtaine trupmena: \(5\frac{3}{10}\).
\end{enumerate}

\subsection{Ryšys su kitomis temomis}

Trupmenos tęsia dalybos idėją iš skyriaus \ref{sec:daugybos-ir-dalybos-pradzia}. Jos reikalingos matuojant ilgį, plotą ir tūrį skyriuose \ref{sec:perimetras-plotas-turis-ir-masto-pradzia} ir \ref{sec:matavimo-vienetai-laikas-greitis-plotas-turis}, taip pat paprastoms tikimybėms skyriuje \ref{sec:paprastos-tikimybes}.

\section{Reiškiniai ir raidiniai žymėjimai}\label{sec:reiskiniai-ir-raidiniai-zymejimai}

\subsection{Apibrėžtys}

\begin{apibreztis}
Reiškinys yra skaičių, raidžių, veiksmų ženklų ir skliaustų užrašas. Pavyzdžiui, \(8+3\cdot4\) ir \(a+5\) yra reiškiniai.
\end{apibreztis}

\begin{apibreztis}
Raidė matematikoje gali žymėti nežinomą arba kintantį skaičių. Pavyzdžiui, jei dėžutėje yra \(x\) pieštukų, tai \(x+3\) reiškia trimis pieštukais daugiau.
\end{apibreztis}

\begin{apibreztis}
Reiškinio reikšmė yra skaičius, gaunamas atlikus veiksmus. Jei \(a=7\), tai reiškinio \(a+5\) reikšmė yra \(12\).
\end{apibreztis}

\subsection{Teoremos ir savybės}

\begin{savybe}
Ta pati situacija gali būti aprašyta žodžiais, piešiniu, lentele, skaitiniu reiškiniu arba raidiniu reiškiniu.
\end{savybe}

\begin{savybe}
Jei į tą patį reiškinį vietoje raidės įrašome skirtingus skaičius, dažniausiai gauname skirtingas reikšmes.
\end{savybe}

\begin{savybe}
Skliaustai reiškinyje parodo, kad pirmiausia reikia atlikti veiksmus jų viduje.
\end{savybe}

\subsection{Taisyklės ir algoritmai}

\textbf{Reiškinio reikšmės radimas.}
\begin{enumerate}[label=\arabic*.]
\item Perskaityk, kokia raidės reikšmė duota.
\item Vietoje raidės įrašyk duotą skaičių.
\item Atlik veiksmus pagal veiksmų tvarką.
\item Parašyk atsakymą su situacijos vienetais, jei jų yra.
\end{enumerate}

\textbf{Situacijos užrašymas reiškiniu.}
\begin{enumerate}[label=\arabic*.]
\item Rask nežinomą arba kintantį kiekį.
\item Pažymėk jį raide.
\item Pagal žodžius parink veiksmą: padidinti -- \(+\), sumažinti -- \(-\), kartų daugiau -- \(\cdot\), padalyti į lygias dalis -- \(:\).
\end{enumerate}

\subsection{Iliustruojantys pavyzdžiai}

\textbf{1 pavyzdys.} Raskime reiškinio \(a+18\) reikšmę, kai \(a=27\).

Vietoje \(a\) įrašome \(27\):
\[
a+18=27+18=45.
\]
\textbf{Atsakymas:} \(45\).

\textbf{2 pavyzdys.} Užrašykime reiškinį: mokinys turėjo \(x\) lipdukų, o draugas jam davė dar \(12\) lipdukų.

Pradinis lipdukų skaičius yra \(x\). Pridedame \(12\), todėl gauname
\[
x+12.
\]
\textbf{Atsakymas:} \(x+12\).

\textbf{3 pavyzdys.} Raskime reiškinio \(5\cdot(n-3)\) reikšmę, kai \(n=11\).

Įrašome \(11\):
\[
5\cdot(n-3)=5\cdot(11-3).
\]
Pirma skliaustai:
\[
11-3=8.
\]
Tada daugyba:
\[
5\cdot8=40.
\]
\textbf{Atsakymas:} \(40\).

\subsection{Dažnos klaidos ir savikontrolė}

\textbf{Dažnos klaidos.}
\begin{enumerate}[label=\arabic*.]
\item Vietoje raidės įrašomas skaičius, bet pamirštami skliaustai.
\item Žodžiai ``\(5\) mažiau negu \(x\)'' klaidingai užrašomi \(5-x\), nors reikia \(x-5\).
\item Reiškinys supainiojamas su lygtimi: reiškinyje nebūtinai yra lygybės ženklas.
\end{enumerate}

\textbf{Pasitikrinkite.}
\begin{enumerate}[label=\arabic*.]
\item Raskite \(b-14\), kai \(b=52\).
\item Užrašykite reiškinį: \(m\) obuolių padalijama po lygiai \(4\) vaikams.
\item Raskite \(3\cdot(k+2)\), kai \(k=8\).
\end{enumerate}

\subsection{Ryšys su kitomis temomis}

Raidiniai reiškiniai remiasi veiksmų tvarka iš skyriaus \ref{sec:aritmetiniai-veiksmai-ir-savybes}. Jie padeda spręsti lygtis skyriuje \ref{sec:paprastos-lygtys-ir-nelygybes}, užrašyti lentelių taisykles skyriuje \ref{sec:sekos-lenteles-ir-koordinates} ir pasiruošti vėlesnei algebrai skyriuje \ref{sec:algebra-lygtys-ir-raidiniai-reiskiniai}.

\section{Paprastos lygtys ir nelygybės}\label{sec:paprastos-lygtys-ir-nelygybes}

\subsection{Apibrėžtys}

\begin{apibreztis}
Lygtis yra lygybė su nežinomuoju, pavyzdžiui \(x+7=15\).
\end{apibreztis}

\begin{apibreztis}
Lygties sprendinys yra skaičius, kurį įrašius vietoje nežinomojo lygybė tampa teisinga.
\end{apibreztis}

\begin{apibreztis}
Nelygybė yra užrašas su ženklais \(<\), \(>\), \(\leq\), \(\geq\). Pavyzdžiui, \(x<9\) reiškia, kad \(x\) yra mažesnis už \(9\).
\end{apibreztis}

\subsection{Teoremos ir savybės}

\begin{savybe}
Jei prie abiejų lygties pusių pridedame tą patį skaičių arba iš abiejų pusių atimame tą patį skaičių, lygties pusiausvyra išlieka.
\end{savybe}

\begin{savybe}
Sudėties lygtis tikrinama atimtimi: jei \(x+6=20\), tai \(x=20-6\).
\end{savybe}

\begin{savybe}
Nelygybė gali turėti daugiau negu vieną sprendinį. Pavyzdžiui, jei ieškome natūraliųjų skaičių, nelygybei \(x<4\) tinka \(1\), \(2\), \(3\).
\end{savybe}

\subsection{Taisyklės ir algoritmai}

\textbf{Lygtys su sudėtimi.}
Nežinomą dėmenį randame iš sumos atimdami žinomą dėmenį.
\[
x+18=43,\quad x=43-18.
\]

\textbf{Lygtys su atimtimi.}
\begin{enumerate}[label=\arabic*.]
\item Jei nežinomas skaičius sumažintas: \(x-5=12\), tada \(x=12+5\).
\item Jei nežinomas atimamas skaičius: \(20-x=8\), tada \(x=20-8\).
\end{enumerate}

\textbf{Nelygybių sprendimas bandant.}
Pasirink skaičius iš nurodytos srities, įrašyk vietoje nežinomojo ir palik tik tuos, kurie padaro nelygybę teisingą.

\subsection{Iliustruojantys pavyzdžiai}

\textbf{1 pavyzdys.} Išspręskime lygtį \(x+18=43\).

Nežinomas dėmuo yra \(x\). Jį randame iš sumos atimdami žinomą dėmenį:
\[
x=43-18=25.
\]
Tikriname:
\[
25+18=43.
\]
\textbf{Atsakymas:} \(x=25\).

\textbf{2 pavyzdys.} Išspręskime lygtį \(y-9=16\).

Nežinomas skaičius sumažintas \(9\) ir gauta \(16\). Vadinasi, prie \(16\) reikia pridėti \(9\):
\[
y=16+9=25.
\]
Tikriname:
\[
25-9=16.
\]
\textbf{Atsakymas:} \(y=25\).

\textbf{3 pavyzdys.} Raskime natūraliuosius skaičius, mažesnius už \(6\), kurie tinka nelygybei \(2\cdot x<9\).

Bandome:
\[
2\cdot1=2<9,\quad 2\cdot2=4<9,\quad 2\cdot3=6<9,\quad 2\cdot4=8<9,\quad 2\cdot5=10>9.
\]
Tinka \(1,2,3,4\).
\textbf{Atsakymas:} \(x=1,2,3,4\).

\subsection{Dažnos klaidos ir savikontrolė}

\textbf{Dažnos klaidos.}
\begin{enumerate}[label=\arabic*.]
\item Sprendinys randamas, bet nepatikrinamas įrašant į pradinę lygtį.
\item Lygtyse \(20-x=8\) klaidingai skaičiuojama \(x=8-20\).
\item Nelygybei pateikiamas tik vienas sprendinys, nors jų gali būti keli.
\end{enumerate}

\textbf{Pasitikrinkite.}
\begin{enumerate}[label=\arabic*.]
\item Išspręskite: \(x+24=70\).
\item Išspręskite: \(m-13=29\).
\item Raskite natūraliuosius \(x<8\), kuriems \(3\cdot x<18\).
\end{enumerate}

\subsection{Ryšys su kitomis temomis}

Lygtys ir nelygybės tęsia nežinomojo idėją iš skyriaus \ref{sec:lygybes-nelygybes-ir-nezinomasis}. Jos naudoja reiškinius iš skyriaus \ref{sec:reiskiniai-ir-raidiniai-zymejimai} ir vėliau plėtojamos skyriuje \ref{sec:lygtys-ir-nelygybes-su-vienu-nezinomuoju}.

\section{Sekos, lentelės ir koordinatės}\label{sec:sekos-lenteles-ir-koordinates}

\subsection{Apibrėžtys}

\begin{apibreztis}
Seka yra skaičių arba objektų eilė, sudaryta pagal tam tikrą taisyklę.
\end{apibreztis}

\begin{apibreztis}
Lentelė yra duomenų arba skaičiavimo rezultatų užrašymo būdas eilutėmis ir stulpeliais.
\end{apibreztis}

\begin{apibreztis}
Koordinatės nusako taško vietą tinklelyje. Porą \((3;2)\) skaitome taip: nuo pradžios einame \(3\) langelius į dešinę ir \(2\) langelius aukštyn.
\end{apibreztis}

\subsection{Teoremos ir savybės}

\begin{savybe}
Sekoje gali būti pastovus žingsnis. Sekoje \(5,10,15,20,\ldots\) žingsnis yra \(5\).
\end{savybe}

\begin{savybe}
Lentelėje ta pati taisyklė taikoma kiekvienai eilutei. Jei taisyklė yra \(y=3\cdot x\), kiekvienas \(x\) dauginamas iš \(3\).
\end{savybe}

\begin{savybe}
Koordinačių tvarka svarbi: taškas \((3;2)\) yra kitoje vietoje negu taškas \((2;3)\).
\end{savybe}

\subsection{Taisyklės ir algoritmai}

\textbf{Sekos pratęsimas.}
\begin{enumerate}[label=\arabic*.]
\item Palygink gretimus narius.
\item Nustatyk, ar jie didėja, mažėja, kartojasi, dauginami ar keičiasi kitaip.
\item Taikyk rastą taisyklę kitam nariui.
\item Patikrink taisyklę su keliais ankstesniais nariais.
\end{enumerate}

\textbf{Lentelės pildymas.}
\begin{enumerate}[label=\arabic*.]
\item Perskaityk taisyklę.
\item Kiekvieną įėjimo skaičių įrašyk į taisyklę.
\item Apskaičiuok išėjimo skaičių.
\item Patikrink, ar visose eilutėse taikyta ta pati taisyklė.
\end{enumerate}

\textbf{Taško žymėjimas tinklelyje.}
Pirmoji koordinatė rodo judėjimą į dešinę, antroji -- judėjimą aukštyn.

\subsection{Iliustruojantys pavyzdžiai}

\textbf{1 pavyzdys.} Pratęskime seką \(12,17,22,27,\ldots\).

Randame žingsnį:
\[
17-12=5,\quad 22-17=5,\quad 27-22=5.
\]
Kitas narys:
\[
27+5=32.
\]
\textbf{Atsakymas:} \(32\).

\textbf{2 pavyzdys.} Užpildykime lentelę, jei \(y=2\cdot x+1\), o \(x=1,2,3\).

\[
\begin{array}{c|c}
x & y\\ \hline
1 & 2\cdot1+1=3\\
2 & 2\cdot2+1=5\\
3 & 2\cdot3+1=7
\end{array}
\]
\textbf{Atsakymas:} \(3,5,7\).

\textbf{3 pavyzdys.} Kuriame taške atsidursime, jei nuo pradžios eisime \(4\) langelius į dešinę ir \(3\) langelius aukštyn?

Pirmoji koordinatė yra \(4\), antroji \(3\), todėl taškas yra
\[
(4;3).
\]
\textbf{Atsakymas:} \((4;3)\).

\subsection{Dažnos klaidos ir savikontrolė}

\textbf{Dažnos klaidos.}
\begin{enumerate}[label=\arabic*.]
\item Sprendžiant seką remiamasi tik paskutiniais dviem nariais ir nepatikrinami ankstesni.
\item Lentelėje vienoms eilutėms taikoma viena taisyklė, kitoms -- kita.
\item Koordinatėse supainiojama pirmoji ir antroji reikšmės.
\end{enumerate}

\textbf{Pasitikrinkite.}
\begin{enumerate}[label=\arabic*.]
\item Pratęskite seką: \(40,36,32,28,\ldots\).
\item Užpildykite lentelę, kai \(y=4\cdot x-2\), o \(x=2,3,4\).
\item Aprašykite, kaip pažymėti tašką \((5;1)\).
\end{enumerate}

\subsection{Ryšys su kitomis temomis}

Sekos tęsia dėsningumų mokymąsi iš skyriaus \ref{sec:desningumai-ir-sekos}. Lentelės siejasi su reiškiniais \ref{sec:reiskiniai-ir-raidiniai-zymejimai} ir duomenų vaizdavimu \ref{sec:duomenu-lenteles-ir-diagramos}. Koordinatės plėtojamos skyriuje \ref{sec:koordinaciu-planai-postumiai-ir-posukiai}.

\section{Kampai, tiesės ir daugiakampiai}\label{sec:kampai-tieses-ir-daugiakampiai}

\subsection{Apibrėžtys}

\begin{apibreztis}
Tiesė yra be galo į abi puses besitęsianti tiesi linija. Atkarpa yra tiesios linijos dalis, turinti du galus.
\end{apibreztis}

\begin{apibreztis}
Spindulys yra tiesios linijos dalis, turinti pradžią ir besitęsianti viena kryptimi.
\end{apibreztis}

\begin{apibreztis}
Kampas susidaro, kai du spinduliai prasideda tame pačiame taške. Tas taškas vadinamas kampo viršūne.
\end{apibreztis}

\begin{apibreztis}
Daugiakampis yra uždara figūra, sudaryta iš atkarpų. Trikampis turi \(3\) kraštines, keturkampis turi \(4\) kraštines, penkiakampis turi \(5\) kraštines.
\end{apibreztis}

\subsection{Teoremos ir savybės}

\begin{savybe}
Statusis kampas yra toks kampas, koks susidaro kvadrato arba stačiakampio kampe.
\end{savybe}

\begin{savybe}
Smailusis kampas yra mažesnis už statųjį, o bukasis kampas yra didesnis už statųjį.
\end{savybe}

\begin{savybe}
Daugiakampio kraštinių skaičius lygus jo viršūnių skaičiui.
\end{savybe}

\begin{savybe}
Kvadratas yra keturkampis, kurio visos kraštinės lygios ir visi kampai statieji. Stačiakampis yra keturkampis, kurio visi kampai statieji, bet gretimos kraštinės gali būti skirtingo ilgio.
\end{savybe}

\subsection{Taisyklės ir algoritmai}

\textbf{Kampo rūšies nustatymas.}
\begin{enumerate}[label=\arabic*.]
\item Palygink kampą su stačiuoju kampu, pavyzdžiui, langelio arba liniuotės kampu.
\item Jei kampas toks pat, jis statusis.
\item Jei kampas mažesnis, jis smailusis.
\item Jei kampas didesnis, jis bukasis.
\end{enumerate}

\textbf{Daugiakampio atpažinimas.}
\begin{enumerate}[label=\arabic*.]
\item Patikrink, ar figūra uždara.
\item Patikrink, ar ji sudaryta iš atkarpų.
\item Suskaičiuok kraštines ir viršūnes.
\item Pagal kraštinių skaičių pavadink figūrą.
\end{enumerate}

\subsection{Iliustruojantys pavyzdžiai}

\textbf{1 pavyzdys.} Figūra turi \(6\) kraštines ir \(6\) viršūnes. Kaip ji vadinama?

Daugiakampis, turintis \(6\) kraštines, vadinamas šešiakampiu.
\textbf{Atsakymas:} šešiakampis.

\textbf{2 pavyzdys.} Ar kiekvienas kvadratas yra stačiakampis?

Stačiakampis turi turėti \(4\) stačiuosius kampus. Kvadratas taip pat turi \(4\) stačiuosius kampus, todėl kvadratas yra stačiakampis. Tačiau ne kiekvienas stačiakampis yra kvadratas, nes stačiakampio kraštinės nebūtinai visos lygios.
\textbf{Atsakymas:} taip, kiekvienas kvadratas yra stačiakampis.

\textbf{3 pavyzdys.} Atkarpos ilgis \(AB=6\) cm, o atkarpos \(CD\) ilgis \(2\) cm didesnis. Koks \(CD\) ilgis?

\[
6+2=8.
\]
\textbf{Atsakymas:} \(8\) cm.

\subsection{Dažnos klaidos ir savikontrolė}

\textbf{Dažnos klaidos.}
\begin{enumerate}[label=\arabic*.]
\item Atvira laužtė laikoma daugiakampiu, nors daugiakampis turi būti uždaras.
\item Kvadratas nelaikomas stačiakampiu, nors jis turi visas stačiakampio savybes.
\item Kampas vertinamas pagal kraštinių ilgį, nors kampo dydį lemia spindulių išsiskyrimas.
\end{enumerate}

\textbf{Pasitikrinkite.}
\begin{enumerate}[label=\arabic*.]
\item Kiek kraštinių turi aštuonkampis?
\item Kuo skiriasi atkarpa ir tiesė?
\item Ar keturkampis su vienu stačiuoju kampu būtinai yra stačiakampis?
\end{enumerate}

\subsection{Ryšys su kitomis temomis}

Geometrinės figūros reikalingos perimetrui, plotui ir tūriui skyriuje \ref{sec:perimetras-plotas-turis-ir-masto-pradzia}, transformacijoms skyriuje \ref{sec:simetrija-ir-transformacijos-pradzia} ir koordinačių planams skyriuje \ref{sec:koordinaciu-planai-postumiai-ir-posukiai}.

\section{Perimetras, plotas, tūris ir masto pradžia}\label{sec:perimetras-plotas-turis-ir-masto-pradzia}

\subsection{Apibrėžtys}

\begin{apibreztis}
Perimetras yra figūros visų kraštinių ilgių suma.
\end{apibreztis}

\begin{apibreztis}
Plotas parodo, kiek plokštumos dalies užima figūra. Pradžioje plotą patogu matuoti vienetiniais kvadratėliais.
\end{apibreztis}

\begin{apibreztis}
Tūris parodo, kiek vietos užima erdvinis kūnas. Pradžioje tūrį patogu matuoti vienetiniais kubeliais.
\end{apibreztis}

\begin{apibreztis}
Mastas rodo, kiek kartų brėžinyje pavaizduotas ilgis yra mažesnis arba didesnis už tikrąjį ilgį.
\end{apibreztis}

\subsection{Teoremos ir savybės}

\begin{savybe}
Stačiakampio perimetras apskaičiuojamas sudedant visas kraštines:
\[
P=a+b+a+b=2\cdot a+2\cdot b.
\]
\end{savybe}

\begin{savybe}
Stačiakampio plotas lygus ilgio ir pločio sandaugai:
\[
S=a\cdot b.
\]
\end{savybe}

\begin{savybe}
Stačiakampio gretasienio tūris lygus ilgio, pločio ir aukščio sandaugai:
\[
V=a\cdot b\cdot c.
\]
III--IV klasėje šią mintį galima suprasti kaip vienodų kubelių sluoksnių skaičiavimą.
\end{savybe}

\subsection{Taisyklės ir algoritmai}

\textbf{Perimetro radimas.}
\begin{enumerate}[label=\arabic*.]
\item Užrašyk visų kraštinių ilgius tuo pačiu matavimo vienetu.
\item Sudėk kraštinių ilgius.
\item Atsakyme parašyk ilgio vienetą.
\end{enumerate}

\textbf{Ploto radimas langeliais.}
\begin{enumerate}[label=\arabic*.]
\item Suskaičiuok, kiek vienetinių kvadratėlių telpa vienoje eilėje.
\item Suskaičiuok, kiek yra eilių.
\item Padaugink šiuos skaičius.
\end{enumerate}

\textbf{Masto skaitymas.}
Jei mastas \(1:100\), tai \(1\) cm brėžinyje atitinka \(100\) cm tikrovėje.

\subsection{Iliustruojantys pavyzdžiai}

\textbf{1 pavyzdys.} Stačiakampio ilgis \(8\) cm, plotis \(5\) cm. Raskime perimetrą.

\[
P=8+5+8+5=26\ \text{cm}.
\]
\textbf{Atsakymas:} \(26\) cm.

\textbf{2 pavyzdys.} Raskime to paties stačiakampio plotą.

\[
S=8\cdot5=40\ \text{cm}^2.
\]
\textbf{Atsakymas:} \(40\ \text{cm}^2\).

\textbf{3 pavyzdys.} Dėžė sudėta iš \(3\) sluoksnių. Kiekviename sluoksnyje yra \(4\cdot5=20\) kubelių. Kiek kubelių dėžėje?

\[
20\cdot3=60.
\]
\textbf{Atsakymas:} \(60\) kubelių.

\textbf{4 pavyzdys.} Plane \(1\) cm atitinka \(2\) m. Kiek metrų atitinka \(6\) cm?

\[
6\cdot2=12.
\]
\textbf{Atsakymas:} \(12\) m.

\subsection{Dažnos klaidos ir savikontrolė}

\textbf{Dažnos klaidos.}
\begin{enumerate}[label=\arabic*.]
\item Perimetras ir plotas supainiojami: perimetras matuojamas ilgio vienetais, plotas -- kvadratiniais vienetais.
\item Sudedami skirtingi vienetai, pavyzdžiui, centimetrai ir metrai, jų nepakeitus.
\item Skaičiuojant tūrį pamirštamas aukštis arba sluoksnių skaičius.
\end{enumerate}

\textbf{Pasitikrinkite.}
\begin{enumerate}[label=\arabic*.]
\item Raskite kvadrato, kurio kraštinė \(7\) cm, perimetrą.
\item Raskite stačiakampio, kurio kraštinės \(9\) cm ir \(4\) cm, plotą.
\item Jei plane \(1\) cm atitinka \(5\) m, kiek metrų atitinka \(3\) cm?
\end{enumerate}

\subsection{Ryšys su kitomis temomis}

Perimetras ir plotas remiasi daugiakampiais iš skyriaus \ref{sec:kampai-tieses-ir-daugiakampiai}, daugyba iš skyriaus \ref{sec:aritmetiniai-veiksmai-ir-savybes} ir matavimo vienetais iš skyriaus \ref{sec:matavimo-vienetai-laikas-greitis-plotas-turis}.

\section{Simetrija ir transformacijos pradžia}\label{sec:simetrija-ir-transformacijos-pradzia}

\subsection{Apibrėžtys}

\begin{apibreztis}
Simetrijos ašis yra tiesė, kuri figūrą padalija į dvi sutampančias dalis.
\end{apibreztis}

\begin{apibreztis}
Atspindys yra figūros vaizdas kitoje simetrijos ašies pusėje taip, tarsi figūra atsispindėtų veidrodyje.
\end{apibreztis}

\begin{apibreztis}
Postūmis yra figūros perkėlimas nurodyta kryptimi ir nurodytu atstumu nekeičiant jos formos ir dydžio.
\end{apibreztis}

\begin{apibreztis}
Posūkis yra figūros pasukimas apie pasirinktą tašką.
\end{apibreztis}

\subsection{Teoremos ir savybės}

\begin{savybe}
Atspindint figūrą, jos forma ir dydis nesikeičia.
\end{savybe}

\begin{savybe}
Atliekant postūmį, visi figūros taškai pajuda ta pačia kryptimi ir tokiu pačiu atstumu.
\end{savybe}

\begin{savybe}
Jei figūra turi simetrijos ašį, vienoje ašies pusėje esanti dalis sutampa su kita dalimi užlenkus per ašį.
\end{savybe}

\subsection{Taisyklės ir algoritmai}

\textbf{Simetrijos ašies tikrinimas.}
\begin{enumerate}[label=\arabic*.]
\item Įsivaizduok, kad figūra užlenkiama per pasirinktą tiesę.
\item Patikrink, ar abi pusės sutampa.
\item Jei sutampa, tiesė yra simetrijos ašis.
\end{enumerate}

\textbf{Atspindžio braižymas tinklelyje.}
\begin{enumerate}[label=\arabic*.]
\item Pasirink figūros viršūnę.
\item Suskaičiuok, kiek langelių ji nutolusi nuo ašies.
\item Kitoje ašies pusėje pažymėk tašką tokiu pačiu atstumu.
\item Tą patį atlik visoms viršūnėms ir sujunk jas atkarpomis.
\end{enumerate}

\subsection{Iliustruojantys pavyzdžiai}

\textbf{1 pavyzdys.} Ar kvadratas turi simetrijos ašių?

Kvadratą galima perlenkti per vidurį vertikaliai, horizontaliai ir per abi įstrižaines. Visais atvejais pusės sutampa.
\textbf{Atsakymas:} kvadratas turi \(4\) simetrijos ašis.

\textbf{2 pavyzdys.} Taškas yra \(3\) langeliai į kairę nuo vertikalios simetrijos ašies. Kur bus jo atspindys?

Atspindys bus \(3\) langeliai į dešinę nuo tos pačios ašies.
\textbf{Atsakymas:} \(3\) langeliai į dešinę nuo ašies.

\textbf{3 pavyzdys.} Figūra pastumta \(4\) langeliais į dešinę ir \(1\) langeliu aukštyn. Ar jos plotas pasikeitė?

Postūmis nekeičia figūros dydžio, todėl plotas nepasikeitė.
\textbf{Atsakymas:} ne, plotas nepasikeitė.

\subsection{Dažnos klaidos ir savikontrolė}

\textbf{Dažnos klaidos.}
\begin{enumerate}[label=\arabic*.]
\item Atspindžio taškas pažymimas ne tokiu pačiu atstumu nuo ašies.
\item Postūmio metu pakeičiamas figūros dydis.
\item Simetrijos ašimi palaikoma bet kuri figūrą kertanti tiesė.
\end{enumerate}

\textbf{Pasitikrinkite.}
\begin{enumerate}[label=\arabic*.]
\item Kiek simetrijos ašių turi stačiakampis, kuris nėra kvadratas?
\item Taškas yra \(2\) langeliai virš horizontalios ašies. Kur jo atspindys?
\item Ar posūkis keičia figūros kraštinių ilgius?
\end{enumerate}

\subsection{Ryšys su kitomis temomis}

Transformacijos remiasi figūrų pažinimu iš skyriaus \ref{sec:kampai-tieses-ir-daugiakampiai}, koordinatėmis iš skyriaus \ref{sec:sekos-lenteles-ir-koordinates} ir toliau plėtojamos skyriuje \ref{sec:koordinaciu-planai-postumiai-ir-posukiai}.

\section{Duomenų lentelės ir diagramos}\label{sec:duomenu-lenteles-ir-diagramos}

\subsection{Apibrėžtys}

\begin{apibreztis}
Duomenys yra surinkta informacija, pavyzdžiui, kiek mokinių pasirinko kiekvieną mėgstamą vaisių.
\end{apibreztis}

\begin{apibreztis}
Lentelė yra duomenų užrašymo būdas eilutėmis ir stulpeliais.
\end{apibreztis}

\begin{apibreztis}
Diagrama yra duomenų vaizdavimas piešiniais, stulpeliais, taškais arba kitais ženklais.
\end{apibreztis}

\begin{apibreztis}
Skalė nurodo, kiek reiškia vienas langelis, padala arba ženklas diagramoje.
\end{apibreztis}

\subsection{Teoremos ir savybės}

\begin{savybe}
Stulpelinėje diagramoje aukštesnis stulpelis rodo didesnį kiekį, jei visiems stulpeliams taikoma ta pati skalė.
\end{savybe}

\begin{savybe}
Duomenis lengviau palyginti, kai jie tvarkingai suskirstyti į grupes ir pateikti lentele arba diagrama.
\end{savybe}

\begin{savybe}
Bendras kiekis gaunamas sudedant visų grupių kiekius.
\end{savybe}

\subsection{Taisyklės ir algoritmai}

\textbf{Lentelės sudarymas.}
\begin{enumerate}[label=\arabic*.]
\item Pirmame stulpelyje užrašyk grupes arba požymius.
\item Antrame stulpelyje užrašyk kiekius.
\item Patikrink, ar kiekvienas duomuo įrašytas tinkamoje grupėje.
\end{enumerate}

\textbf{Diagramos skaitymas.}
\begin{enumerate}[label=\arabic*.]
\item Perskaityk diagramos pavadinimą.
\item Pažiūrėk, ką reiškia stulpeliai arba ženklai.
\item Patikrink skalę.
\item Pagal skalę perskaityk kiekį.
\end{enumerate}

\textbf{Klausimo atsakymas pagal duomenis.}
Pirmiausia rask reikalingus skaičius lentelėje arba diagramoje, tada atlik tinkamą veiksmą: sudėk, atimk, palygink arba padaugink pagal skalę.

\subsection{Iliustruojantys pavyzdžiai}

\textbf{1 pavyzdys.} Lentelėje parodyta, kiek mokinių pasirinko vaisius:
\[
\begin{array}{c|c}
\text{Vaisius} & \text{Mokinių skaičius}\\ \hline
\text{Obuolys} & 9\\
\text{Bananas} & 6\\
\text{Kriaušė} & 5
\end{array}
\]
Kiek mokinių atsakė iš viso?

\[
9+6+5=20.
\]
\textbf{Atsakymas:} \(20\) mokinių.

\textbf{2 pavyzdys.} Kiek daugiau mokinių pasirinko obuolį nei bananą?

\[
9-6=3.
\]
\textbf{Atsakymas:} \(3\) mokiniais daugiau.

\textbf{3 pavyzdys.} Diagramos skalė: vienas langelis reiškia \(2\) mokinius. Stulpelis yra \(7\) langelių aukščio. Kiek mokinių jis rodo?

\[
7\cdot2=14.
\]
\textbf{Atsakymas:} \(14\) mokinių.

\subsection{Dažnos klaidos ir savikontrolė}

\textbf{Dažnos klaidos.}
\begin{enumerate}[label=\arabic*.]
\item Skaitoma stulpelio aukštis langeliais, bet nepaisoma skalės.
\item Sumaišomi klausimai ``kiek iš viso'' ir ``kiek daugiau''.
\item Diagramoje lyginami stulpeliai su skirtingomis skalėmis.
\end{enumerate}

\textbf{Pasitikrinkite.}
\begin{enumerate}[label=\arabic*.]
\item Lentelėje yra \(8\), \(12\), \(10\) atsakymų. Kiek atsakymų iš viso?
\item Kiek daugiau yra \(15\) negu \(9\)?
\item Jei vienas paveikslėlis reiškia \(5\) vaikus, kiek vaikų reiškia \(6\) paveikslėliai?
\end{enumerate}

\subsection{Ryšys su kitomis temomis}

Duomenų lentelės tęsia duomenų rinkimą iš skyriaus \ref{sec:duomenu-rinkimas-ir-diagramos}. Jos siejasi su lentelių pildymu skyriuje \ref{sec:sekos-lenteles-ir-koordinates}, veiksmais skyriuje \ref{sec:aritmetiniai-veiksmai-ir-savybes} ir tikimybėmis skyriuje \ref{sec:paprastos-tikimybes}.

\section{Paprastos tikimybės}\label{sec:paprastos-tikimybes}

\subsection{Apibrėžtys}

\begin{apibreztis}
Įvykis yra tai, kas gali atsitikti bandymo metu. Pavyzdžiui, metant kauliuką įvykis gali būti ``iškrito \(6\)''.
\end{apibreztis}

\begin{apibreztis}
Būtinas įvykis visada atsitinka, negalimas įvykis niekada neatsitinka, o galimas įvykis gali atsitikti arba neatsitikti.
\end{apibreztis}

\begin{apibreztis}
Paprasta tikimybė nusako, kiek palankių rezultatų yra tarp visų vienodai galimų rezultatų.
\end{apibreztis}

\subsection{Teoremos ir savybės}

\begin{savybe}
Jei maišelyje daugiau raudonų kortelių negu mėlynų, ištraukti raudoną kortelę yra labiau tikėtina negu mėlyną.
\end{savybe}

\begin{savybe}
Jei visi rezultatai vienodai galimi, tikimybę galima užrašyti trupmena:
\[
\frac{\text{palankių rezultatų skaičius}}{\text{visų rezultatų skaičius}}.
\]
\end{savybe}

\begin{savybe}
Tikimybė negali būti mažesnė už \(0\) ir negali būti didesnė už \(1\). Pradinėse klasėse ją suprantame kaip galimybių dalį.
\end{savybe}

\subsection{Taisyklės ir algoritmai}

\textbf{Įvykio apibūdinimas žodžiais.}
\begin{enumerate}[label=\arabic*.]
\item Paklausk, ar įvykis tikrai atsitiks.
\item Jei tikrai atsitiks, jis būtinas.
\item Jei negali atsitikti, jis negalimas.
\item Jei gali atsitikti, bet nebūtinai, jis galimas.
\end{enumerate}

\textbf{Paprastos tikimybės radimas.}
\begin{enumerate}[label=\arabic*.]
\item Suskaičiuok visus vienodai galimus rezultatus.
\item Suskaičiuok palankius rezultatus.
\item Sudaryk trupmeną.
\item Jei reikia, palygink tikimybes pagal palankių rezultatų skaičių.
\end{enumerate}

\subsection{Iliustruojantys pavyzdžiai}

\textbf{1 pavyzdys.} Maišelyje yra \(3\) raudoni ir \(2\) mėlyni rutuliukai. Kiek iš viso rutuliukų?

\[
3+2=5.
\]
\textbf{Atsakymas:} \(5\) rutuliukai.

\textbf{2 pavyzdys.} Kokia tikimybė ištraukti raudoną rutuliuką?

Visų rutuliukų yra \(5\). Palankūs rezultatai yra \(3\) raudoni rutuliukai. Todėl tikimybė
\[
\frac{3}{5}.
\]
\textbf{Atsakymas:} \(\frac{3}{5}\).

\textbf{3 pavyzdys.} Metamas įprastas žaidimo kauliukas. Ar įvykis ``iškris skaičius \(8\)'' yra galimas?

Ant įprasto kauliuko yra skaičiai \(1,2,3,4,5,6\). Skaičiaus \(8\) nėra, todėl toks įvykis negali atsitikti.
\textbf{Atsakymas:} negalimas įvykis.

\subsection{Dažnos klaidos ir savikontrolė}

\textbf{Dažnos klaidos.}
\begin{enumerate}[label=\arabic*.]
\item Palankūs rezultatai supainiojami su visais rezultatais.
\item Teigiama, kad retas įvykis yra negalimas, nors jis tik mažiau tikėtinas.
\item Tikimybė užrašoma didesnė už \(1\), pavyzdžiui \(\frac{7}{5}\), kai kalbama apie vieną bandymą su \(5\) vienodai galimomis baigtimis.
\end{enumerate}

\textbf{Pasitikrinkite.}
\begin{enumerate}[label=\arabic*.]
\item Maišelyje yra \(4\) žalios ir \(6\) geltonos kortelės. Kiek kortelių iš viso?
\item Kokia tikimybė ištraukti žalią kortelę?
\item Ar metant kauliuką įvykis ``iškris skaičius mažesnis už \(7\)'' yra būtinas?
\end{enumerate}

\subsection{Ryšys su kitomis temomis}

Tikimybės žodžiai remiasi ankstesne tema \ref{sec:tikimybes-zodziai-ir-ivykiai}. Tikimybės naudoja trupmenas iš skyriaus \ref{sec:trupmenos-ir-desimtaines-trupmenos-pradzia} ir duomenų skaitymą iš skyriaus \ref{sec:duomenu-lenteles-ir-diagramos}.

\section{Skaičiai nuo 0 iki 10 000}\label{sec:skaiciai-nuo-0-iki-10000}

\subsection{Apibrėžtys}

\begin{apibreztis}
Skaičiai nuo \(0\) iki \(10\,000\) apima nulį ir visus natūraliuosius skaičius iki dešimties tūkstančių. Šioje temoje ypač svarbu mokėti skaičius skaityti, rašyti, lyginti, apvalinti ir skaidyti skyrių suma.
\end{apibreztis}

\begin{apibreztis}
Skaičių tiesė yra tiesė, kurioje skaičiai pažymėti didėjimo tvarka. Ji padeda matyti, kuris skaičius didesnis, mažesnis, arčiau nurodyto skaičiaus.
\end{apibreztis}

\subsection{Teoremos ir savybės}

\begin{savybe}
Keturiženklį skaičių sudaro tūkstančiai, šimtai, dešimtys ir vienetai:
\[
7\,408=7\,000+400+8.
\]
\end{savybe}

\begin{savybe}
Skaičius \(10\,000\) yra vienu didesnis už \(9\,999\):
\[
9\,999+1=10\,000.
\]
\end{savybe}

\begin{savybe}
Kuo skaičius skaičių tiesėje yra dešiniau, tuo jis didesnis.
\end{savybe}

\subsection{Taisyklės ir algoritmai}

\textbf{Keturiženklio skaičiaus nagrinėjimas.}
\begin{enumerate}[label=\arabic*.]
\item Nustatyk tūkstančių skaitmenį.
\item Nustatyk šimtų, dešimčių ir vienetų skaitmenis.
\item Užrašyk skyrių sumą.
\item Perskaityk skaičių žodžiais.
\end{enumerate}

\textbf{Apvalinimas iki dešimčių, šimtų ir tūkstančių.}
Taikoma ta pati apvalinimo taisyklė kaip skyriuje \ref{sec:naturalieji-skaiciai-iki-milijono}: žiūrime į dešinėje esantį skaitmenį ir sprendžiame, ar apvalinamą skaitmenį didinti.

\subsection{Iliustruojantys pavyzdžiai}

\textbf{1 pavyzdys.} Užrašykime \(6\,305\) skyrių suma.

\[
6\,305=6\,000+300+5.
\]
\textbf{Atsakymas:} \(6\,000+300+5\).

\textbf{2 pavyzdys.} Palyginkime \(4\,908\) ir \(4\,980\).

Tūkstančiai vienodi: \(4=4\). Šimtai vienodi: \(9=9\). Dešimtys skiriasi: \(0<8\), todėl
\[
4\,908<4\,980.
\]
\textbf{Atsakymas:} \(4\,908<4\,980\).

\textbf{3 pavyzdys.} Suapvalinkime \(7\,642\) iki šimtų.

Šimtų skaitmuo yra \(6\), dešinėje yra dešimčių skaitmuo \(4\). Kadangi \(4<5\), šimtų skaitmens nedidiname:
\[
7\,642\approx7\,600.
\]
\textbf{Atsakymas:} \(7\,600\).

\subsection{Dažnos klaidos ir savikontrolė}

\textbf{Dažnos klaidos.}
\begin{enumerate}[label=\arabic*.]
\item Skaičius \(6\,305\) klaidingai skaitomas kaip \(635\), nes nepastebimas nulis dešimčių vietoje.
\item Apvalinant iki šimtų žiūrima į vienetų, o ne į dešimčių skaitmenį.
\item Skaičių tiesėje didesnis skaičius pažymimas kairiau.
\end{enumerate}

\textbf{Pasitikrinkite.}
\begin{enumerate}[label=\arabic*.]
\item Užrašykite skyrių suma: \(9\,072\).
\item Įrašykite ženklą: \(5\,601\ \square\ 5\,610\).
\item Suapvalinkite \(8\,549\) iki tūkstančių.
\end{enumerate}

\subsection{Ryšys su kitomis temomis}

Skaičiai iki \(10\,000\) yra tiltas nuo ankstesnių skaičių iki \(1000\) prie skaičių iki milijono skyriuje \ref{sec:naturalieji-skaiciai-iki-milijono}. Jie dažnai naudojami pinigų, matavimo ir duomenų uždaviniuose.

\section{Finansiniai skaičiavimai: kainos, pokyčiai ir sprendimai}\label{sec:finansiniai-skaiciavimai-kainos-pokyciai-sprendimai}

\subsection{Apibrėžtys}

\begin{apibreztis}
Kaina yra pinigų suma, kurią reikia sumokėti už prekę arba paslaugą.
\end{apibreztis}

\begin{apibreztis}
Išlaidos yra pinigai, kuriuos sumokame. Santaupos yra pinigai, kuriuos pasiliekame arba kaupiame vėlesniam tikslui.
\end{apibreztis}

\begin{apibreztis}
Grąža yra pinigų suma, kurią gauname atgal, kai sumokame daugiau, negu kainuoja pirkinys.
\end{apibreztis}

\begin{apibreztis}
Nuolaida yra kainos sumažinimas. III--IV klasėje dažniausiai nagrinėjamos paprastos nuolaidos eurais arba centais.
\end{apibreztis}

\subsection{Teoremos ir savybės}

\begin{savybe}
Bendra pirkinių kaina gaunama sudedant visų pirkinių kainas.
\end{savybe}

\begin{savybe}
Grąža apskaičiuojama iš sumokėtos sumos atimant pirkinio kainą.
\end{savybe}

\begin{savybe}
Pinigų sumas galima stambinti ir smulkinti: \(1\) euras yra \(100\) centų.
\end{savybe}

\subsection{Taisyklės ir algoritmai}

\textbf{Pirkinių kainos radimas.}
\begin{enumerate}[label=\arabic*.]
\item Užrašyk kiekvienos prekės kainą.
\item Jei perkama keli vienodi daiktai, kainą padaugink iš kiekio.
\item Sudėk visas kainas.
\item Patikrink, ar vienetai tie patys: eurai su eurais, centai su centais.
\end{enumerate}

\textbf{Grąžos radimas.}
\begin{enumerate}[label=\arabic*.]
\item Užrašyk, kiek pinigų sumokėta.
\item Užrašyk pirkinio kainą.
\item Atimk kainą iš sumokėtos sumos.
\end{enumerate}

\textbf{Sprendimo pasirinkimas.}
Jei reikia pasirinkti, ką pirkti, palygink kainą, kiekį, poreikį ir turimą pinigų sumą. Pigiausias pasirinkimas ne visada yra tinkamiausias, jei jis neatitinka uždavinio sąlygos.

\subsection{Iliustruojantys pavyzdžiai}

\textbf{1 pavyzdys.} Sąsiuvinis kainuoja \(2\) eurus, pieštukas \(70\) centų. Kiek kainuoja \(2\) sąsiuviniai ir \(1\) pieštukas?

\[
2\cdot2=4\ \text{eurai}.
\]
Pridedame \(70\) centų:
\[
4\ \text{Eur}+70\ \text{ct}=4\ \text{Eur}\ 70\ \text{ct}.
\]
\textbf{Atsakymas:} \(4\) Eur \(70\) ct.

\textbf{2 pavyzdys.} Pirkėjas sumokėjo \(10\) Eur už prekes, kurios kainavo \(7\) Eur \(35\) ct. Kokia grąža?

\[
10\ \text{Eur}=9\ \text{Eur}\ 100\ \text{ct}.
\]
\[
9\ \text{Eur}\ 100\ \text{ct}-7\ \text{Eur}\ 35\ \text{ct}=2\ \text{Eur}\ 65\ \text{ct}.
\]
\textbf{Atsakymas:} \(2\) Eur \(65\) ct.

\textbf{3 pavyzdys.} Kuprinė kainavo \(28\) Eur. Jai pritaikyta \(5\) Eur nuolaida. Kokia nauja kaina?

\[
28-5=23.
\]
\textbf{Atsakymas:} \(23\) Eur.

\subsection{Dažnos klaidos ir savikontrolė}

\textbf{Dažnos klaidos.}
\begin{enumerate}[label=\arabic*.]
\item Eurai ir centai sudedami kaip paprasti skaičiai: \(4\) Eur \(70\) ct nėra \(4,70\) ct.
\item Grąža skaičiuojama atvirkščiai: iš kainos atimama sumokėta suma.
\item Perkant kelias vienodas prekes pamirštama dauginti iš kiekio.
\end{enumerate}

\textbf{Pasitikrinkite.}
\begin{enumerate}[label=\arabic*.]
\item Kiek kainuoja \(3\) prekės po \(4\) Eur?
\item Kokia grąža iš \(20\) Eur, jei pirkinys kainuoja \(13\) Eur \(80\) ct?
\item Prekė kainavo \(15\) Eur, nuolaida \(3\) Eur. Kokia nauja kaina?
\end{enumerate}

\subsection{Ryšys su kitomis temomis}

Finansiniai skaičiavimai taiko sudėtį, atimtį, daugybą ir dalybą iš skyriaus \ref{sec:aritmetiniai-veiksmai-ir-savybes}, skaičius iki \(10\,000\) iš skyriaus \ref{sec:skaiciai-nuo-0-iki-10000} ir tekstinių uždavinių modeliavimą iš skyriaus \ref{sec:reiskiniai-ir-raidiniai-zymejimai}.

\section{Algoritmai: pasirinkimas ir kartojimas}\label{sec:algoritmai-pasirinkimas-ir-kartojimas}

\subsection{Apibrėžtys}

\begin{apibreztis}
Algoritmas yra aiškių žingsnių seka, kurią atlikę išsprendžiame užduotį arba pasiekiame tikslą.
\end{apibreztis}

\begin{apibreztis}
Pasirinkimo komanda nurodo, ką daryti skirtingais atvejais. Ji dažnai turi formą: jei sąlyga teisinga, atlik vieną veiksmą, priešingu atveju atlik kitą.
\end{apibreztis}

\begin{apibreztis}
Kartojimo komanda nurodo tą patį veiksmą atlikti kelis kartus arba tol, kol sąlyga pasikeičia.
\end{apibreztis}

\subsection{Teoremos ir savybės}

\begin{savybe}
Aiškus algoritmas turi baigtinį žingsnių skaičių arba aiškią kartojimo pabaigos sąlygą.
\end{savybe}

\begin{savybe}
Tas pats algoritmas, taikomas tiems patiems duomenims, turi duoti tą patį rezultatą, jei žingsniai atliekami vienodai.
\end{savybe}

\begin{savybe}
Algoritmo klaidą dažnai galima rasti tikrinant žingsnius po vieną.
\end{savybe}

\subsection{Taisyklės ir algoritmai}

\textbf{Algoritmo sudarymas.}
\begin{enumerate}[label=\arabic*.]
\item Aiškiai įvardyk tikslą.
\item Išvardyk veiksmus tokia tvarka, kokia juos reikia atlikti.
\item Nurodyk pasirinkimus: kada ką daryti.
\item Nurodyk kartojimus: kiek kartų arba iki kada kartoti.
\item Išbandyk algoritmą su paprastu pavyzdžiu.
\end{enumerate}

\textbf{Pasirinkimo pavyzdys.}
Jei skaičius didesnis už \(100\), jį apvalink iki šimtų; priešingu atveju palik nepakeistą.

\textbf{Kartojimo pavyzdys.}
Kartok: pridėk \(5\), kol pasieksi arba viršysi \(40\).

\subsection{Iliustruojantys pavyzdžiai}

\textbf{1 pavyzdys.} Algoritmas: pradėk nuo \(12\), tris kartus pridėk \(4\). Koks rezultatas?

\[
12+4=16,\quad 16+4=20,\quad 20+4=24.
\]
\textbf{Atsakymas:} \(24\).

\textbf{2 pavyzdys.} Pasirinkimo taisyklė: jei skaičius lyginis, padalyk iš \(2\), jei nelyginis, pridėk \(1\). Ką gausime iš \(18\)?

Skaičius \(18\) lyginis, todėl dalijame:
\[
18:2=9.
\]
\textbf{Atsakymas:} \(9\).

\textbf{3 pavyzdys.} Kartojame: pradėk nuo \(3\), daugink iš \(2\), kol rezultatas bus didesnis už \(20\).

\[
3\cdot2=6,\quad 6\cdot2=12,\quad 12\cdot2=24.
\]
Skaičius \(24\) jau didesnis už \(20\), todėl sustojame.
\textbf{Atsakymas:} \(24\).

\subsection{Dažnos klaidos ir savikontrolė}

\textbf{Dažnos klaidos.}
\begin{enumerate}[label=\arabic*.]
\item Kartojimas neturi pabaigos sąlygos.
\item Sąlyga perskaitoma atvirkščiai: atliekamas ``kitaip'' veiksmas, nors sąlyga teisinga.
\item Žingsniai surašomi netvarkinga tvarka, todėl rezultatas pasikeičia.
\end{enumerate}

\textbf{Pasitikrinkite.}
\begin{enumerate}[label=\arabic*.]
\item Pradėkite nuo \(5\) ir keturis kartus pridėkite \(3\). Ką gausite?
\item Jei skaičius mažesnis už \(50\), pridėkite \(10\); kitaip atimkite \(10\). Ką gausite iš \(47\)?
\item Kodėl algoritme svarbu nurodyti, kada kartojimas baigiasi?
\end{enumerate}

\subsection{Ryšys su kitomis temomis}

Algoritmai padeda aiškiai atlikti stulpelinę daugybą ir dalybą skyriuje \ref{sec:daugyba-ir-dalyba-stulpeliu}, spręsti lygtis skyriuje \ref{sec:paprastos-lygtys-ir-nelygybes}, pildyti lenteles skyriuje \ref{sec:sekos-lenteles-ir-koordinates} ir tikrinti uždavinių sprendimus.

\section{Algebra: lygtys ir raidiniai reiškiniai}\label{sec:algebra-lygtys-ir-raidiniai-reiskiniai}

\subsection{Apibrėžtys}

\begin{apibreztis}
Algebra pradinėse klasėse prasideda tada, kai skaičius pradedame žymėti raidėmis, sudarome reiškinius, lygtis ir ieškome nežinomų reikšmių.
\end{apibreztis}

\begin{apibreztis}
Raidinis reiškinys yra reiškinys, kuriame yra bent viena raidė, pavyzdžiui \(2\cdot a+5\).
\end{apibreztis}

\begin{apibreztis}
Paprastas matematinis modelis yra užrašas, schema, lentelė arba lygtis, kuri aprašo realią situaciją.
\end{apibreztis}

\subsection{Teoremos ir savybės}

\begin{savybe}
Raidė gali reikšti skirtingus skaičius skirtingose situacijose, todėl visada reikia žiūrėti, kas pasakyta sąlygoje.
\end{savybe}

\begin{savybe}
Lygtį galima suprasti kaip pusiausvyrą: abi pusės turi būti lygios.
\end{savybe}

\begin{savybe}
Tą patį tekstinį uždavinį dažnai galima spręsti keliais būdais: veiksmų seka, lentele, schema arba lygtimi.
\end{savybe}

\subsection{Taisyklės ir algoritmai}

\textbf{Tekstinio uždavinio modeliavimas lygtimi.}
\begin{enumerate}[label=\arabic*.]
\item Perskaityk klausimą ir nuspręsk, kas nežinoma.
\item Nežinomą dydį pažymėk raide.
\item Pagal sąlygą sudaryk lygybę.
\item Išspręsk lygtį.
\item Patikrink, ar atsakymas tinka pradinei situacijai.
\end{enumerate}

\textbf{Raidinio reiškinio reikšmės lentelė.}
Jei reikia palyginti kelias reikšmes, patogu sudaryti lentelę: pirmame stulpelyje rašyti raidės reikšmes, antrame -- reiškinio reikšmes.

\subsection{Iliustruojantys pavyzdžiai}

\textbf{1 pavyzdys.} Ieva turėjo kelis lipdukus. Gavusi dar \(9\), ji turi \(31\). Kiek lipdukų Ieva turėjo iš pradžių?

Nežinomą pradinį skaičių pažymime \(x\). Tada
\[
x+9=31.
\]
\[
x=31-9=22.
\]
\textbf{Atsakymas:} \(22\) lipdukus.

\textbf{2 pavyzdys.} Raskime \(2\cdot a+5\), kai \(a=6\).

\[
2\cdot a+5=2\cdot6+5=12+5=17.
\]
\textbf{Atsakymas:} \(17\).

\textbf{3 pavyzdys.} Užpildykime lentelę reiškiniui \(p-4\), kai \(p=10,15,20\).

\[
\begin{array}{c|c}
p & p-4\\ \hline
10 & 6\\
15 & 11\\
20 & 16
\end{array}
\]
\textbf{Atsakymas:} \(6,11,16\).

\subsection{Dažnos klaidos ir savikontrolė}

\textbf{Dažnos klaidos.}
\begin{enumerate}[label=\arabic*.]
\item Nežinomasis pasirenkamas ne pagal klausimą, todėl lygtis neatitinka situacijos.
\item Patikrinama tik skaičiavimo klaida, bet ne tai, ar atsakymas prasmingas.
\item Raidė laikoma tik vienu konkrečiu skaičiumi, nors kitoje eilutėje ji gali turėti kitą duotą reikšmę.
\end{enumerate}

\textbf{Pasitikrinkite.}
\begin{enumerate}[label=\arabic*.]
\item Sudarykite lygtį: skaičių padidinus \(12\), gauta \(40\).
\item Raskite \(3\cdot b-2\), kai \(b=7\).
\item Sudarykite lentelę reiškiniui \(x+6\), kai \(x=4,8,12\).
\end{enumerate}

\subsection{Ryšys su kitomis temomis}

Šis skyrius apibendrina reiškinius iš skyriaus \ref{sec:reiskiniai-ir-raidiniai-zymejimai} ir lygtis iš skyriaus \ref{sec:paprastos-lygtys-ir-nelygybes}. Jis ruošia aukštesnių klasių algebros temoms, kuriose lygtys ir reiškiniai bus sudėtingesni.

\section{Matavimo vienetai: laikas, greitis, plotas ir tūris}\label{sec:matavimo-vienetai-laikas-greitis-plotas-turis}

\subsection{Apibrėžtys}

\begin{apibreztis}
Matavimo vienetas yra susitartas dydis, kuriuo matuojame ilgį, masę, laiką, plotą, tūrį ar kitą dydį.
\end{apibreztis}

\begin{apibreztis}
Laikas matuojamas sekundėmis, minutėmis, valandomis, paromis, savaitėmis, mėnesiais ir metais.
\end{apibreztis}

\begin{apibreztis}
Greitis parodo, kokį kelią objektas nueina arba nuvažiuoja per vieną laiko vienetą. Pradžioje greitį suprantame per paprastus pavyzdžius: jei per \(1\) valandą nuvažiuojama \(60\) km, greitis yra \(60\) km per valandą.
\end{apibreztis}

\begin{apibreztis}
Plotas matuojamas kvadratiniais vienetais, pavyzdžiui \(\text{cm}^2\), o tūris -- kubiniais vienetais, pavyzdžiui \(\text{cm}^3\), arba litrais.
\end{apibreztis}

\subsection{Teoremos ir savybės}

\begin{savybe}
\[
1\ \text{val.}=60\ \text{min.},\quad 1\ \text{min.}=60\ \text{s}.
\]
\end{savybe}

\begin{savybe}
\[
1\ \text{m}=100\ \text{cm},\quad 1\ \text{km}=1000\ \text{m}.
\]
\end{savybe}

\begin{savybe}
Jei judama pastoviu greičiu, kelias randamas greitį dauginant iš laiko:
\[
\text{kelias}=\text{greitis}\cdot\text{laikas}.
\]
III--IV klasėje ši taisyklė taikoma tik paprastoms gyvenimiškoms situacijoms.
\end{savybe}

\subsection{Taisyklės ir algoritmai}

\textbf{Vienetų keitimas.}
\begin{enumerate}[label=\arabic*.]
\item Nuspręsk, ar keiti į mažesnius, ar į didesnius vienetus.
\item Keičiant į mažesnius vienetus, dažniausiai dauginama.
\item Keičiant į didesnius vienetus, dažniausiai dalijama.
\item Patikrink, ar atsakymas prasmingas.
\end{enumerate}

\textbf{Laiko trukmės radimas.}
\begin{enumerate}[label=\arabic*.]
\item Užrašyk pradžios ir pabaigos laiką.
\item Skaičiuok iki artimiausios valandos, tada pilnas valandas, tada likusias minutes.
\item Sudėk gautas dalis.
\end{enumerate}

\subsection{Iliustruojantys pavyzdžiai}

\textbf{1 pavyzdys.} Kiek minučių yra \(3\) valandose?

\[
3\cdot60=180.
\]
\textbf{Atsakymas:} \(180\) min.

\textbf{2 pavyzdys.} Kiek centimetrų yra \(4\) metruose \(25\) centimetruose?

\[
4\ \text{m}=400\ \text{cm},
\]
\[
400\ \text{cm}+25\ \text{cm}=425\ \text{cm}.
\]
\textbf{Atsakymas:} \(425\) cm.

\textbf{3 pavyzdys.} Dviratininkas per \(2\) valandas nuvažiavo po \(12\) km kiekvieną valandą. Kokį kelią jis nuvažiavo?

\[
12\cdot2=24.
\]
\textbf{Atsakymas:} \(24\) km.

\textbf{4 pavyzdys.} Stačiakampio ilgis \(6\) cm, plotis \(3\) cm. Raskime plotą.

\[
S=6\cdot3=18\ \text{cm}^2.
\]
\textbf{Atsakymas:} \(18\ \text{cm}^2\).

\subsection{Dažnos klaidos ir savikontrolė}

\textbf{Dažnos klaidos.}
\begin{enumerate}[label=\arabic*.]
\item Keičiant metrus į centimetrus dalijama iš \(100\), nors reikia dauginti.
\item Laiko uždaviniuose skaičiuojama kaip dešimtainėje sistemoje, pavyzdžiui, \(1\) val. \(80\) min. vietoje \(2\) val. \(20\) min.
\item Ploto ir tūrio vienetai rašomi kaip paprasti ilgio vienetai.
\end{enumerate}

\textbf{Pasitikrinkite.}
\begin{enumerate}[label=\arabic*.]
\item Kiek sekundžių yra \(5\) minutėse?
\item Kiek metrų yra \(3\) km?
\item Per \(3\) valandas nueinama po \(4\) km kas valandą. Koks kelias?
\end{enumerate}

\subsection{Ryšys su kitomis temomis}

Matavimo vienetai remiasi skaičiavimais iš skyriaus \ref{sec:aritmetiniai-veiksmai-ir-savybes}. Plotas ir tūris siejasi su skyriumi \ref{sec:perimetras-plotas-turis-ir-masto-pradzia}, o greičio lentelės -- su skyriumi \ref{sec:sekos-lenteles-ir-koordinates}.

\section{Koordinačių planai, postūmiai ir posūkiai}\label{sec:koordinaciu-planai-postumiai-ir-posukiai}

\subsection{Apibrėžtys}

\begin{apibreztis}
Koordinačių planas yra tinklelis, kuriame taškų vieta nusakoma skaičių poromis.
\end{apibreztis}

\begin{apibreztis}
Koordinačių pradžia yra taškas, nuo kurio pradedame skaičiuoti judėjimą į dešinę ir aukštyn.
\end{apibreztis}

\begin{apibreztis}
Postūmis koordinačių plane aprašo, kiek langelių figūra perkelta į dešinę arba kairę ir kiek langelių aukštyn arba žemyn.
\end{apibreztis}

\begin{apibreztis}
Posūkis koordinačių plane yra figūros pasukimas apie nurodytą tašką. Pradinėse klasėse dažniausiai nagrinėjami posūkiai pagal langelių tinklelį ir paprastus kampus.
\end{apibreztis}

\subsection{Teoremos ir savybės}

\begin{savybe}
Postūmis nekeičia figūros formos, dydžio, kraštinių ilgių ir kampų.
\end{savybe}

\begin{savybe}
Posūkis nekeičia figūros formos ir dydžio, tik pakeičia jos padėtį.
\end{savybe}

\begin{savybe}
Taško koordinatės \((a;b)\) reiškia, kad pirmiausia skaičiuojame horizontalią kryptį, paskui vertikalią kryptį.
\end{savybe}

\subsection{Taisyklės ir algoritmai}

\textbf{Taško radimas plane.}
\begin{enumerate}[label=\arabic*.]
\item Pradėk nuo koordinačių pradžios.
\item Pagal pirmą koordinatę eik į dešinę.
\item Pagal antrą koordinatę eik aukštyn.
\item Pažymėk tašką.
\end{enumerate}

\textbf{Postūmio aprašymas.}
\begin{enumerate}[label=\arabic*.]
\item Pasirink vieną figūros viršūnę.
\item Suskaičiuok, kiek langelių ji pajudėjo į dešinę arba kairę.
\item Suskaičiuok, kiek langelių ji pajudėjo aukštyn arba žemyn.
\item Patikrink kitą viršūnę, ar postūmis toks pats.
\end{enumerate}

\subsection{Iliustruojantys pavyzdžiai}

\textbf{1 pavyzdys.} Taškas \(A=(2;3)\) pastumtas \(4\) langeliais į dešinę. Kokios naujos koordinatės?

Pirmoji koordinatė padidėja \(4\), antroji nesikeičia:
\[
2+4=6.
\]
\textbf{Atsakymas:} \(A'=(6;3)\).

\textbf{2 pavyzdys.} Taškas \(B=(5;6)\) pastumtas \(2\) langeliais žemyn. Kokios naujos koordinatės?

Pirmoji koordinatė nesikeičia, antroji sumažėja \(2\):
\[
6-2=4.
\]
\textbf{Atsakymas:} \(B'=(5;4)\).

\textbf{3 pavyzdys.} Figūra pasukta apie vieną viršūnę. Ar jos kraštinių ilgiai pasikeitė?

Posūkis nekeičia ilgių, todėl kraštinių ilgiai liko tokie patys.
\textbf{Atsakymas:} ne, nepasikeitė.

\subsection{Dažnos klaidos ir savikontrolė}

\textbf{Dažnos klaidos.}
\begin{enumerate}[label=\arabic*.]
\item Pirma judama aukštyn, nors pirmoji koordinatė rodo judėjimą į dešinę.
\item Pastumiama tik viena figūros viršūnė, o kitos paliekamos vietoje.
\item Posūkis supainiojamas su atspindžiu.
\end{enumerate}

\textbf{Pasitikrinkite.}
\begin{enumerate}[label=\arabic*.]
\item Tašką \((1;4)\) pastumkite \(3\) langeliais į dešinę. Kokios koordinatės?
\item Tašką \((6;2)\) pastumkite \(1\) langeliu aukštyn. Kokios koordinatės?
\item Ar postūmis gali pakeisti kvadrato plotą?
\end{enumerate}

\subsection{Ryšys su kitomis temomis}

Koordinačių planai tęsia koordinates iš skyriaus \ref{sec:sekos-lenteles-ir-koordinates} ir transformacijas iš skyriaus \ref{sec:simetrija-ir-transformacijos-pradzia}. Jie taip pat padeda skaityti planus, brėžinius ir paprastus žemėlapius.
